\documentclass[11pt,oneside]{report}

\usepackage{semantic_motion_report}
\usepackage[numbers,sort&compress]{natbib}

\ProjectTitle{Semantic-Motion}
\ProjectSubtitle{A Multi-view Engineering Pipeline for Perception, Fact-preserving Language Augmentation and Fail-closed Multimodal Refinement}
\AuthorName{Your Name}
\StudentID{Your Student ID}
\SupervisorName{Your Supervisor}
\SubmissionDate{Month 2026}

\begin{document}

\makeprojecttitle
\frontmatterstyle

\begin{declaration}
State originality, academic integrity, use of generative AI where required,
and individual contributions for a group project.
\end{declaration}

\begin{projectabstract}
Keep the abstract below 200 words. State the data problem, the four-module
pipeline, the evaluated outputs, the strongest quantitative result, and the
main limitation. Do not describe planned work as completed work.
\end{projectabstract}

\tableofcontents
\listoffigures
\listoftables

\mainmatterstyle

\chapter{Introduction}
\section{Project Context: From Motion Data to VLA-ready Multimodal Data}
\section{Engineering Problem and Motivation}
\section{Project Aim and Engineering Deliverables}
\section{System Requirements}
\section{Scope and Boundaries}
\section{Contributions}
\section{Report Organisation}

\chapter{Technical Background and Related Systems}
\section{Vision-Language-Action Models}
\section{Robot Learning Benchmarks}
\section{Vision-Language Models for Robotics}
\section{Video-to-Task Generation}
\section{Demonstration Data Quality}
\section{Engineering Gap}

\chapter{System Requirements and Overall Design}
\section{Inputs, Outputs and Use Cases}
\section{Design Goals}
\section{Overall Architecture}
\section{Module Responsibilities and Boundaries}
\section{Data Models and Interface Contracts}
\subsection{Paired Multi-view Video Contract}
\subsection{Shared Timebase and Motion Trajectory Contract}
\section{Observability and Visualisation Design}
\section{End-to-End Workflow}

\chapter{Design and Implementation of the Four Modules}
\section{The Rendering Module}
\subsection{Motion Data Import and Scene Initialisation}
\subsection{Scene Context Generation}
\subsection{Multi-view Video Synthesis}
\subsection{3D Motion Trajectory Export}
\subsection{Rendering Outputs and Manifest}

\section{The Perception Module}
\subsection{Paired-view Input and Synchronous Frame Sampling}
\subsection{Overlapping Macro Windows and Boundary Voting}
\subsection{VLM-based Structured Scene Understanding}
\subsection{Macro-Intent Recognition}
\subsection{Dense Micro-Instruction Parsing}
\subsection{Body, Contact, Posture and Trajectory Evidence}
\subsection{Temporal Task Boundary Detection}
\subsection{Segment-level Temporal Aggregation}
\subsection{Static Perception Benchmark}

\section{The Fact-preserving LLM Augmentation Module}
\subsection{Input and Interface}
\subsection{Instruction Reconstruction}
\subsection{Augmentation Strategies}
\subsection{Deterministic Fact Preservation and Hallucination Control}
\subsection{Output and Provenance}

\section{The Refinement Module}
\subsection{Sample Construction}
\subsection{Motion Data Normalisation}
\subsection{Motion Quality Scoring}
\subsection{Semantic Consistency Verification}
\subsection{Deterministic Cross-modal Alignment Checks}
\subsection{Fail-closed Refinement Decision and Explainable Outputs}

\chapter{System Integration, Demonstration and Evaluation}
\section{Research Questions}
\section{Integration Environment and Technology Stack}
\section{Module Integration and Interface Verification}
\section{End-to-End Demonstration and Visualisation}
\section{Evaluation Setup}
\section{Perception Module Evaluation}
\section{Pinned Video2Tasks Full-pipeline Comparison}
\section{Refinement and Alignment Evaluation}
\subsection{Human Gold Annotation Protocol}
\subsection{Temporal Boundary and Segment Metrics}
\subsection{Keep/Drop and Confidence Calibration}
\subsection{Controlled Motion Corruptions}
\section{Cross-module Case Studies}
\section{Component and Configuration Analysis}
\section{Engineering Acceptance Summary}

\chapter{Engineering Discussion}
\section{Interpretation of System Results}
\section{Engineering Strengths}
\section{Failure Mode Analysis}
\section{Limitations}
\section{Engineering Trade-offs}
\section{Evaluation Validity}
\section{Implications for VLA Data Pipelines}

\chapter{Conclusion and Future Work}
\section{Conclusion}
\noindent Report only automatically verified or independently evaluated results.
Mark API, Blender, and pending-human evidence explicitly; do not infer completion.
\section{Future Work}

\bibliographystyle{plainnat}
\bibliography{references}

\appendix
\chapter{VLM System Prompts}
\chapter{Semantic Consistency Prompt}
\chapter{Data Schemas}
\chapter{Configuration}
\chapter{Additional Results}
\chapter{User Guide}

\end{document}
